\chapter{Phân hệ Đặt hàng}

Phân hệ này gồm \textbf{14 đường dẫn} gom thành \textbf{6 màn} người dùng nhìn thấy, nằm trong
mô-đun \rt{wujia_portal_sale} (tệp \rt{controllers/portal.py}). Đây là phân hệ phức tạp nhất của
portal: nó vừa hiển thị hàng hoá, vừa giữ giỏ, vừa sinh ra đơn hàng thật trong ERP.

\textbf{Bốn nguyên tắc đang chạy}, cần nắm trước khi đọc từng màn:

\begin{enumerate}[leftmargin=*]
  \item \textbf{Giỏ hàng là của CỬA HÀNG, không phải của người dùng.} Mỗi cửa hàng có đúng một
        giỏ (\rt{wujia.portal.cart}, ràng buộc duy nhất trên cửa hàng). Chủ, quản lý và nhân
        viên cùng cửa hàng thao tác trên cùng một giỏ; ai sửa sau thì giá trị của người đó
        thắng. Người khác đang mở màn giỏ sẽ thấy thay đổi \textbf{ngay lập tức} (bản tin
        realtime theo kênh của cửa hàng).
  \item \textbf{Giá luôn do máy chủ tính}, không tin số do trình duyệt gửi lên. Giá lấy theo
        \emph{bảng giá} gắn với khách hàng của cửa hàng, rồi qua bộ tính thuế của Odoo. Cùng một
        hàm dựng giá được dùng cho lưới sản phẩm, trang chi tiết, giỏ và màn kết quả --- nên
        năm chỗ này không bao giờ lệch số.
  \item \textbf{Đơn gửi từ portal luôn ở trạng thái Nháp}, hệ thống không tự xác nhận. Ngô Gia
        xử lý tiếp trong ERP.
  \item \textbf{Một cửa hàng chỉ có một đơn nháp từ portal.} Gửi đơn mới thì đơn nháp/đã gửi cũ
        của chính cửa hàng đó bị huỷ trong cùng một lượt.
\end{enumerate}

\section{Màn Danh sách sản phẩm}
\rt{/portal/order} · \rt{/portal/order/results}
\medskip

\mvao

Vào từ mục \emph{Đặt hàng} trên thanh điều hướng (cột trái ở máy tính, thanh dưới ở điện thoại).
Đường dẫn thứ hai \rt{/portal/order/results} \textbf{không phải màn riêng}: khi người dùng gõ
tìm kiếm hoặc bấm danh mục, trình duyệt gọi ngầm đường dẫn này để lấy về đúng phần kết quả và
thay tại chỗ, không tải lại cả trang. Hai đường dẫn dùng chung một bộ dữ liệu, nên mọi luật mô
tả dưới đây đúng cho cả hai.

\mai

\begin{itemize}[leftmargin=*]
  \item Phải \textbf{đăng nhập}. Khách vãng lai bị đẩy về trang đăng nhập.
  \item Nếu tài khoản \textbf{không thuộc cửa hàng nào} còn hiệu lực, màn hiện trạng thái rỗng
        ``chưa có cửa hàng'' --- không có sản phẩm, không có giỏ.
  \item \textbf{Cửa hàng đang chọn} được lấy từ \emph{cookie} \rt{wujia_active_franchise_id}
        (lưu 30 ngày trên máy người dùng), và chỉ được chấp nhận nếu cửa hàng đó nằm trong danh
        sách cửa hàng mà tài khoản có thành viên còn hiệu lực. Ai chỉ thuộc \textbf{một} cửa
        hàng thì hệ thống tự chọn, không hỏi.
  \item \textbf{Vai trò không ảnh hưởng màn này}: Chủ, Quản lý hay Nhân viên đều xem và thêm vào
        giỏ như nhau. Phân biệt vai trò chỉ xuất hiện ở các phân hệ khác (công nợ, thành viên).
  \item Nếu \textbf{chưa chọn được cửa hàng} (tài khoản có nhiều cửa hàng mà chưa bấm chọn),
        màn \textbf{vẫn hiện đủ sản phẩm} nhưng: giá rơi về \emph{giá niêm yết} của sản phẩm
        thay vì giá theo bảng giá cửa hàng, giỏ hàng trống, và có dải cảnh báo
        ``Vui lòng chọn cửa hàng trước khi đặt hàng''. Đây là điểm dễ hiểu nhầm --- xem mục 7.
  \item Dữ liệu sản phẩm được đọc \textbf{bằng quyền hệ thống} (bỏ qua phân quyền của người
        dùng portal). Việc giới hạn nằm hoàn toàn ở điều kiện lọc bên dưới, không dựa vào
        quy tắc quyền của Odoo.
\end{itemize}

\mhien

\begin{hienbang}
Lưới sản phẩm &
  \rt{product.product} &
  Bật cờ \emph{Hiện trên portal} (\rt{is_public_portal = true}) \textbf{và} chưa lưu trữ
  (\rt{active = true}). Có từ khoá thì thêm: tên \textbf{hoặc} mã sản phẩm chứa từ khoá. Có
  chọn danh mục thì thêm: đúng danh mục đó. &
  Theo tên A$\rightarrow$Z; 24 sản phẩm/trang &
  \uat{5}/8 sản phẩm đang bật cờ \\ \addlinespace
Bộ lọc danh mục &
  \rt{wujia.product.category} &
  Chỉ những danh mục \textbf{đang có ít nhất một sản phẩm hiện trên portal}. Danh mục đã lưu
  trữ tự động không hiện. &
  Theo \emph{Thứ tự} rồi tên &
  \uat{2} danh mục (Trà, Topping) \\ \addlinespace
Giá từng sản phẩm &
  Bảng giá của khách hàng gắn với cửa hàng &
  Lấy bảng giá ở trường \emph{Bảng giá} của khách hàng cửa hàng; chưa có bảng giá (hoặc chưa
  chọn cửa hàng) thì dùng \emph{giá niêm yết} của sản phẩm. Số hiển thị là giá \textbf{đã gồm
  thuế}, tính theo vị thế tài khoán của chính khách hàng đó. &
  --- &
  4/4 cửa hàng đều đã gắn khách hàng \\ \addlinespace
Số trong giỏ trên mỗi thẻ &
  \rt{wujia.portal.cart.line} &
  Dòng thuộc giỏ của cửa hàng đang chọn. &
  --- &
  \uat{6} dòng trong \uat{2} giỏ \\ \addlinespace
Dải khung giờ đặt hàng &
  \rt{wujia.order.window} + tham số chung &
  Theo khu vực của cửa hàng; khu vực chưa cấu hình thì rơi về khung giờ chung. Chi tiết ở
  \textbf{mục 3.7}. &
  --- &
  \uat{1} khung theo khu vực / 2 khu vực \\
\end{hienbang}

\textbf{Trường nào hiện lên thẻ sản phẩm:} ảnh sản phẩm, tên (\rt{name}), tên tiếng Trung
(\rt{name_chinese}, để trống thì không hiện dòng đó), mã (\rt{default_code}), quy cách đóng gói
(\rt{wujia_packaging}, trống thì hiện dấu ``---''), đơn vị tính, giá đã gồm thuế, số lượng tối
thiểu (\rt{min_qty}) dùng làm \textbf{bước tăng giảm}.

\mloc

\begin{itemize}[leftmargin=*]
  \item \textbf{Ô tìm kiếm} --- tìm theo \emph{tên} hoặc \emph{mã sản phẩm}, không phân biệt hoa
        thường, khớp một phần. \textbf{Không} tìm theo tên tiếng Trung và \textbf{không} tìm
        theo quy cách.
  \item \textbf{Danh mục} --- một danh mục tại một thời điểm. Giá trị lạ (không phải số) bị bỏ
        qua, coi như không lọc.
  \item \textbf{Số sản phẩm mỗi trang} --- chỉ nhận \textbf{24, 48, 96}; gửi số khác thì hệ
        thống lặng lẽ dùng 24.
  \item \textbf{Trang} --- số trang vượt quá trang cuối thì tự kéo về trang cuối, không báo lỗi.
        Đổi trang giữ nguyên từ khoá và danh mục đang chọn.
\end{itemize}

\mkiem{Thêm vào giỏ}

Nút này không tải lại trang mà gọi ngầm \rt{/portal/order/cart/add}. Thứ tự kiểm đúng như code
đang chạy:

\begin{kiembang}
1 & Đã chọn được cửa hàng hợp lệ &
    ``Vui lòng chọn cửa hàng trước khi đặt hàng.'' --- hoặc ``Bạn không có quyền thao tác tại
    cửa hàng này'' / ``Tài khoản của bạn hiện không còn hiệu lực tại cửa hàng này'' nếu cookie
    còn trỏ tới cửa hàng cũ \\
2 & Mã sản phẩm gửi lên là số & ``Dữ liệu gửi lên không hợp lệ.'' \\
3 & Sản phẩm còn tồn tại, chưa lưu trữ, còn bật cờ hiện trên portal &
    ``Sản phẩm này hiện không còn được phép đặt hàng.'' \\
4 & Sản phẩm đã cấu hình số lượng tối thiểu (\rt{min_qty > 0}) &
    ``Sản phẩm chưa được cấu hình số lượng đặt tối thiểu. Vui lòng liên hệ Ngô Gia.'' \\
5 & Số lượng gửi lên (nếu có) là \textbf{số nguyên} &
    ``Số lượng của \emph{tên sản phẩm} phải là số nguyên.'' --- 1,5 bị từ chối, \textbf{không}
    làm tròn \\
6 & Số lượng $\geq$ tối thiểu & ``Số lượng tối thiểu của \emph{tên sản phẩm} là \emph{n}.'' \\
7 & Số lượng là \textbf{bội số} của tối thiểu &
    ``Số lượng của \emph{tên sản phẩm} phải tăng theo bước \emph{n}.'' \\
8 & Số lượng $\leq$ tối đa (khi \rt{max_qty > 0}) &
    ``Số lượng tối đa của \emph{tên sản phẩm} là \emph{n}.'' \\
\end{kiembang}

Bấm từ lưới sản phẩm thì trình duyệt \textbf{không gửi số lượng}, hệ thống tự cộng đúng một
bước (bằng số lượng tối thiểu). Các bước 5--8 chỉ áp dụng khi người dùng gõ số tay ở trang chi
tiết. Nếu cộng dồn vượt trần tối đa, hệ thống \textbf{cắt xuống đúng trần} và trả về cảnh báo
chứ không báo lỗi.

\mghi

\begin{itemize}[leftmargin=*]
  \item Tạo giỏ cho cửa hàng nếu chưa có, rồi \textbf{thêm hoặc cộng dồn} dòng
        \rt{wujia.portal.cart.line} (mỗi sản phẩm một dòng duy nhất trong giỏ).
  \item Thao tác cộng dồn chạy thẳng bằng một câu lệnh cơ sở dữ liệu nên hai người cùng bấm một
        lúc \textbf{không mất số của nhau}.
  \item Phát một bản tin realtime tới kênh của cửa hàng để mọi thiết bị khác đang mở giỏ tự cập
        nhật.
  \item Không ghi gì vào đơn hàng --- giỏ và đơn là hai thứ tách rời.
\end{itemize}

\mhoi

\begin{ask}
\begin{enumerate}[leftmargin=*]
  \item \textbf{Sản phẩm chưa gán danh mục vẫn hiện.} Trên UAT có \uat{2}/5 sản phẩm bật cờ
        portal nhưng bỏ trống \emph{Danh mục portal} (\emph{Hồng Trà Sữa size M},
        \emph{Matcha Latte size M}). Chúng hiện ở danh sách chung nhưng \textbf{không bao giờ
        xuất hiện} khi người dùng bấm lọc theo danh mục. Anh muốn (a) bắt buộc phải có danh mục
        mới cho hiện, hay (b) thêm nhóm ``Khác''?
  \item \textbf{Chưa chọn cửa hàng thì giá đang là giá niêm yết.} Người dùng nhiều cửa hàng, khi
        chưa bấm chọn, vẫn thấy giá --- nhưng là giá gốc, không phải giá theo bảng giá của cửa
        hàng. Có nên ẩn hẳn giá cho tới khi chọn cửa hàng không?
  \item \textbf{Tìm kiếm chưa bao gồm tên tiếng Trung.} Màn có hiện \rt{name_chinese} nhưng ô
        tìm kiếm không tìm theo trường đó. Có cần bổ sung không?
  \item \textbf{Trên UAT chưa sản phẩm nào đặt trần tối đa} (\rt{max_qty} đều bằng 0 = không
        giới hạn). Toàn bộ luật chặn ``số lượng tối đa'' vì vậy chưa từng chạy thật khi BA test.
\end{enumerate}
\end{ask}

\section{Màn Chi tiết sản phẩm}
\rt{/portal/order/product/<int:product_id>}
\medskip

\mvao

Bấm vào một thẻ trong lưới sản phẩm.

\mai

Như màn danh sách. Riêng màn này, nếu sản phẩm \textbf{không tồn tại, đã lưu trữ, hoặc đã tắt cờ
hiện trên portal}, người dùng bị đưa về màn danh sách kèm thông báo ``Sản phẩm này hiện không
còn được phép đặt hàng'' --- hệ thống \textbf{cố ý không nói rõ} là sản phẩm không tồn tại hay
người dùng không được xem, để không lộ dữ liệu.

\mhien

\begin{hienbang}
Khối thông tin sản phẩm &
  \rt{product.product} &
  Chính sản phẩm được mở, phải đang \rt{is_public_portal = true} và \rt{active = true} &
  --- & --- \\ \addlinespace
Gợi ý ``sản phẩm liên quan'' &
  \rt{product.product} &
  Cùng \emph{danh mục portal} với sản phẩm đang xem, trừ chính nó, và cũng phải đang hiện trên
  portal. Sản phẩm \textbf{không có danh mục} thì khối này rỗng. &
  Theo tên A$\rightarrow$Z; tối đa 6 &
  2/5 sản phẩm không có danh mục $\Rightarrow$ khối rỗng \\ \addlinespace
Giá &
  Bảng giá của cửa hàng &
  Như màn danh sách --- cùng một hàm tính, nên hai màn luôn cùng số. &
  --- & --- \\
\end{hienbang}

\mloc

Không có bộ lọc. Người dùng chỉ gõ được \textbf{số lượng} trước khi bấm thêm vào giỏ.

\mkiem{Thêm vào giỏ}

Cùng đúng 8 bước kiểm của màn danh sách. Khác một điểm: ở đây trình duyệt \textbf{gửi số lượng
tường minh}, nên các bước 5--8 (số nguyên, tối thiểu, bội số, tối đa) thực sự chạy.

\mghi

Giống màn danh sách.

\mhoi

\begin{ask}
Khối ``sản phẩm liên quan'' đang dựa hoàn toàn vào danh mục portal. Với dữ liệu hiện tại, sản
phẩm không có danh mục sẽ có khối này trống trơn. Anh muốn thay bằng ``sản phẩm hay đặt cùng''
hay để nguyên?
\end{ask}

\section{Màn Giỏ hàng}
\rt{/portal/order/cart} --- kèm 6 đường dẫn thao tác ngầm
\medskip

\mvao

Bấm biểu tượng giỏ trên thanh trên cùng, hoặc nút ``Xem giỏ hàng''. Sáu đường dẫn còn lại không
phải màn, mà là các thao tác trình duyệt gọi ngầm khi người dùng bấm nút:

\begin{itemize}[leftmargin=*,itemsep=1pt]
  \item \rt{/portal/order/cart/add} --- thêm sản phẩm (đã mô tả ở màn danh sách)
  \item \rt{/portal/order/cart/step} --- bấm nút \emph{+} / \emph{--} (tăng giảm đúng một bước)
  \item \rt{/portal/order/cart/update} --- gõ thẳng số lượng vào ô
  \item \rt{/portal/order/cart/remove} --- bấm xoá dòng
  \item \rt{/portal/order/cart/count} --- lấy số hiển thị trên biểu tượng giỏ
  \item \rt{/portal/order/cart/note} --- lưu ghi chú đơn hàng
  \item \rt{/portal/order/cart/fragment} --- lấy lại toàn bộ giỏ khi có người khác vừa sửa
\end{itemize}

\mai

Người dùng \textbf{chưa thuộc cửa hàng nào} bị đưa thẳng về màn danh sách sản phẩm. Ngoài ra như
màn danh sách: giỏ là giỏ chung của cửa hàng đang chọn, ai trong cửa hàng cũng sửa được, không
phân biệt vai trò.

\mhien

\begin{hienbang}
Danh sách dòng &
  \rt{wujia.portal.cart.line} &
  Thuộc giỏ của cửa hàng đang chọn (mỗi cửa hàng đúng một giỏ) &
  Không phân trang &
  \uat{6} dòng / \uat{2} giỏ \\ \addlinespace
Đơn giá từng dòng &
  Bảng giá của cửa hàng &
  Tính lại \textbf{theo số lượng của chính dòng đó} --- bảng giá có bậc số lượng thì giá đổi
  theo. Giỏ \textbf{không lưu giá}, mỗi lần mở là tính lại. &
  --- & --- \\ \addlinespace
Tạm tính · Thuế · Tổng thanh toán &
  Tính tại máy chủ &
  Cộng từ các dòng, làm tròn theo tiền tệ của bảng giá. Trình duyệt không tự cộng. &
  --- & --- \\ \addlinespace
Cảnh báo trên dòng &
  Tính tại máy chủ &
  Mỗi dòng được soi lại ba điều: sản phẩm còn được phép đặt không · đã cấu hình tối thiểu chưa ·
  số lượng có đúng tối thiểu/bội số/tối đa không. &
  --- & --- \\ \addlinespace
Ghi chú đơn hàng &
  \rt{wujia.portal.cart} trường \rt{note} &
  Ghi chú \textbf{dùng chung cả cửa hàng}, tối đa 1000 ký tự. &
  --- & --- \\
\end{hienbang}

\mloc

Không có bộ lọc. Người dùng sửa được số lượng từng dòng và ghi chú chung.

\mkiem{+ / -- / gõ số lượng / xoá}

\begin{kiembang}
1 & Đã chọn được cửa hàng hợp lệ & Thông báo như bảng ở màn danh sách \\
2 & Dòng còn tồn tại và \textbf{thuộc đúng giỏ của cửa hàng đang chọn} &
    \textbf{Không báo lỗi}: coi như dòng đã bị người cùng cửa hàng xoá, trả về giỏ mới nhất và
    màn tự cập nhật \\
3 & Sản phẩm đã cấu hình số lượng tối thiểu &
    ``Sản phẩm chưa cấu hình số lượng tối thiểu.'' \\
4 & Số lượng mới $\geq$ tối thiểu, đúng bội số, $\leq$ tối đa &
    Thông báo tương ứng (tối thiểu / bước / tối đa) \\
\end{kiembang}

Hai hành vi đáng chú ý: \textbf{gõ số lượng bằng 0} hoặc \textbf{bấm ``--'' xuống dưới mức tối
thiểu} đều \textbf{xoá dòng khỏi giỏ}, không báo lỗi. Và nút \emph{+} không cho trình duyệt gửi
số lượng tuỳ ý --- nó chỉ gửi \emph{hướng} tăng hay giảm, máy chủ tự lấy bước --- nên hai người
bấm cùng lúc không ghi đè nhau.

\mghi

\begin{itemize}[leftmargin=*]
  \item Sửa hoặc xoá dòng trong giỏ của cửa hàng; lưu ghi chú vào giỏ.
  \item Mỗi thay đổi phát một bản tin realtime tới kênh của cửa hàng.
  \item Không đụng tới đơn hàng.
\end{itemize}

\mhoi

\begin{ask}
\begin{enumerate}[leftmargin=*]
  \item \textbf{Ghi chú là của cửa hàng, không của người gửi.} Nhân viên A gõ ghi chú, nhân
        viên B gửi đơn thì ghi chú của A đi theo đơn. Anh xác nhận đúng ý nghiệp vụ chứ?
  \item \textbf{Giỏ không có lịch sử.} Khi hai người sửa chồng nhau, hệ thống lấy giá trị của
        người sửa sau và không lưu lại ai đã sửa gì. Có cần nhật ký thao tác giỏ không?
\end{enumerate}
\end{ask}

\section{Gửi đơn hàng}
\rt{/portal/order/submit}
\medskip

\mvao

Nút \emph{Gửi đơn} ở màn giỏ. Đây là đường dẫn \textbf{ghi dữ liệu thật} --- nó tạo đơn hàng
trong ERP, nên có chống giả mạo biểu mẫu (CSRF) và chỉ nhận phương thức POST.

\mai

Chỉ người có thành viên còn hiệu lực tại cửa hàng đang chọn. Vai trò \textbf{không} giới hạn:
nhân viên cũng gửi được đơn.

\mhien

Đường dẫn này \textbf{không hiện màn nào} --- nó chỉ nhận lệnh, ghi dữ liệu rồi chuyển hướng.
Gửi thành công: trên điện thoại sang màn \emph{Đặt hàng thành công} (mục 3.5), trên máy tính sang
thẳng trang chi tiết đơn ở phân hệ \emph{Lịch sử mua hàng}. Gửi hỏng: quay về màn giỏ kèm dải đỏ,
riêng trường hợp ngoài khung giờ trên điện thoại thì sang màn \emph{Không thể gửi đơn} (mục 3.6).

\mloc

\textbf{Không có ô nhập nào.} Toàn bộ nội dung đơn lấy từ giỏ của cửa hàng đang chọn --- trình
duyệt \textbf{không} gửi lên danh sách sản phẩm, số lượng hay giá. Đây là chủ đích: người dùng
không thể can thiệp vào nội dung đơn bằng cách sửa dữ liệu gửi lên.

\mkiem{Gửi đơn}

Đây là chuỗi kiểm dài nhất của portal --- \textbf{11 lớp}, đúng thứ tự chạy:

\begin{kiembang}
1 & Chọn được cửa hàng hợp lệ & Về màn giỏ, dải đỏ theo từng trường hợp (chưa chọn / không có
    quyền / hết hiệu lực) \\
2 & Cửa hàng còn tồn tại & ``Bạn không có quyền thao tác tại cửa hàng này.'' \\
3 & Cửa hàng \textbf{không bị khoá đặt hàng} (\rt{portal_locked}) &
    ``Cửa hàng đang tạm khóa đặt hàng. Vui lòng liên hệ Ngô Gia.'' \\
4 & Giỏ có ít nhất một dòng & ``Giỏ hàng chưa có sản phẩm.'' \\
5 & \textbf{Không có người khác đang gửi đơn cùng lúc} (giỏ bị khoá trong lúc xử lý) &
    ``Giỏ hàng đang được một người dùng khác gửi đơn. Vui lòng thử lại sau.'' \\
6 & \textbf{Đang trong khung giờ đặt hàng} &
    Điện thoại: chuyển sang màn riêng ``Không thể gửi đơn'' có ghi giờ mở lại. Máy tính: về giỏ
    kèm dải ``Hiện ngoài khung giờ đặt hàng.'' \\
7 & Mọi dòng trong giỏ đều hợp lệ (sản phẩm còn bán, số lượng đúng bước) &
    ``Một số sản phẩm trong giỏ không còn được phép đặt'' hoặc ``Số lượng một số sản phẩm chưa
    hợp lệ'' \\
8 & Cửa hàng đã gắn \textbf{khách hàng} &
    ``Cửa hàng chưa được cấu hình khách hàng đặt hàng. Vui lòng liên hệ Ngô Gia.'' \\
9 & Đơn nháp cũ của cửa hàng \textbf{không bị khoá} &
    ``Không thể hoàn tất đơn hàng do đơn nháp cũ chưa được xử lý.'' \\
10 & Tạo được đơn hàng (ERP không từ chối) &
    ``Không thể tạo đơn hàng. Vui lòng thử lại hoặc liên hệ Ngô Gia.'' \\
11 & Huỷ được đơn nháp cũ \textbf{và} dọn được giỏ &
    Hoàn tác toàn bộ --- kể cả đơn vừa tạo --- rồi báo lỗi. Không bao giờ để tồn tại hai đơn
    nháp. \\
\end{kiembang}

Lớp 6 được kiểm \textbf{hai lần}: một lần ở màn, một lần nữa ngay trong ERP lúc tạo đơn. Kể cả
khi ai đó gọi thẳng đường dẫn, đơn ngoài khung giờ vẫn bị chặn.

\mghi

\begin{itemize}[leftmargin=*]
  \item \textbf{Tạo đơn} \rt{sale.order} ở trạng thái \textbf{Nháp}, đánh dấu \rt{is_portal_order},
        gắn cửa hàng, khách hàng của cửa hàng, bảng giá tương ứng, nguồn ``Wujia Portal'',
        người gửi và bản ghi thành viên tại thời điểm gửi.
  \item Ghi chú của giỏ được chép sang đơn ở \textbf{hai trường}: một trường văn bản thuần cho
        trang Lịch sử của portal, một trường ghi chú chuẩn của Odoo.
  \item Mỗi dòng giỏ thành một dòng đơn (sản phẩm, số lượng, đơn vị tính). \textbf{Giá không
        chép sang} --- Odoo tự tính lại theo bảng giá đã gắn, nên giá trên đơn khớp giá người
        dùng vừa thấy.
  \item \textbf{Huỷ} mọi đơn portal cũ của cửa hàng đang ở Nháp/Đã gửi.
  \item \textbf{Dọn sạch} dòng và ghi chú của giỏ, rồi phát bản tin realtime.
  \item Người đứng tên đơn là \textbf{chính tài khoản portal} đã bấm gửi.
\end{itemize}

\mhoi

\begin{ask}
\begin{enumerate}[leftmargin=*]
  \item \textbf{Nhân viên gửi được đơn.} Hiện không phân biệt vai trò ở bước gửi. Nếu nghiệp vụ
        muốn ``chỉ Chủ/Quản lý được gửi đơn'' thì đây là chỗ cần thêm luật.
  \item \textbf{Huỷ đơn nháp cũ là im lặng.} Cửa hàng gửi đơn lần hai trong ngày sẽ mất đơn lần
        một mà không có cảnh báo trước. Có cần hộp xác nhận ``đơn cũ sẽ bị huỷ'' không?
  \item \textbf{Đơn không tự xác nhận.} Hiện Ngô Gia phải vào ERP xác nhận. Đúng như thiết kế
        chứ, hay cần tự xác nhận trong khung giờ?
\end{enumerate}
\end{ask}

\section{Màn Đặt hàng thành công}
\rt{/portal/order/submitted/<int:order_id>}
\medskip

\mvao

Tự chuyển tới sau khi gửi đơn thành công \textbf{trên điện thoại}. Trên máy tính, người dùng
được đưa thẳng sang trang chi tiết đơn ở phân hệ Lịch sử mua hàng --- hai luồng khác nhau có chủ ý.

\mai

Chỉ xem được đơn \textbf{của chính cửa hàng đang chọn} và phải là đơn đặt qua portal. Sai điều
kiện --- kể cả khi gõ tay số đơn của cửa hàng khác --- thì bị đưa về màn danh sách sản phẩm,
không báo gì thêm. Đơn đã huỷ cũng bị chặn ở đây, vì màn này chỉ nói về đơn vừa gửi thành công.

\mhien

\begin{hienbang}
Mã đơn, ngày giờ gửi &
  \rt{sale.order} &
  Đúng đơn vừa tạo; ngày giờ quy về múi giờ người dùng &
  --- &
  \uat{23} đơn đặt từ portal \\ \addlinespace
Tạm tính · Thuế · Tổng &
  \rt{sale.order} &
  Lấy thẳng số tiền của đơn, ký hiệu tiền tệ theo chính đơn đó (không viết cứng ``đ'') &
  --- & --- \\ \addlinespace
Trạng thái đơn &
  \rt{sale.order} trường \rt{state} &
  Dùng \textbf{chung bảng nhãn} với trang Lịch sử mua hàng --- nháp hiển thị là
  ``Chờ xác nhận''. &
  --- &
  \uat{3} đơn portal đang chờ xác nhận \\
\end{hienbang}

\mloc\ Không có ô nhập nào. \quad
\mkiem{Về trang chủ / Xem chi tiết đơn}\ Hai nút trên màn chỉ là liên kết sang màn khác,
không kiểm gì. \quad \mghi\ Không ghi gì --- màn chỉ đọc.

\mhoi

\begin{ask}
Máy tính và điện thoại kết thúc ở \textbf{hai màn khác nhau} sau khi gửi đơn (máy tính sang
trang chi tiết đơn, điện thoại ở màn xác nhận này). Anh xác nhận giữ nguyên chứ?
\end{ask}

\section{Màn Không thể gửi đơn}
\rt{/portal/order/rejected}
\medskip

\mvao

Chỉ dùng cho \textbf{một trường hợp duy nhất}: gửi đơn ngoài khung giờ, trên điện thoại. Mọi lý
do khác đều quay về màn giỏ kèm dải cảnh báo, vì ở đó người dùng sửa được ngay. Nếu ai gõ tay
đường dẫn với lý do lạ, hệ thống đưa về giỏ --- \textbf{không hiển thị lại chuỗi người dùng tự
gõ}, tránh bị lợi dụng để hiện nội dung giả.

\mai\ Như màn giỏ.

\mhien

\begin{hienbang}
Khung giờ mở lại &
  \rt{wujia.order.window} hoặc tham số chung &
  Lấy khung gần nhất sắp mở; nếu bây giờ đã qua giờ mở của mọi khung thì lấy khung đầu tiên của
  \textbf{ngày mai}. Có ghi rõ ``Hôm nay'' hay ngày cụ thể. &
  --- &
  Khu vực Hà Nội chưa cấu hình $\Rightarrow$ dùng khung chung 10:00--04:00 \\
\end{hienbang}

\mloc\ Không có ô nhập nào. \quad
\mkiem{Quay lại giỏ hàng}\ Nút duy nhất trên màn chỉ là liên kết về màn giỏ, không kiểm gì.
\quad \mghi\ Không ghi gì.

\mhoi\ Không có điểm treo.

\section{Khung giờ đặt hàng --- luật dùng chung}
\label{sec:khung-gio}

Đây là luật bị hỏi nhiều nhất, nên tách riêng. Hệ thống quyết định ``có đang trong giờ đặt hàng
không'' theo \textbf{ba bậc}:

\begin{enumerate}[leftmargin=*]
  \item \textbf{Công tắc tổng}. Tham số \rt{portal_order_time_limit_enabled}. Tắt thì đặt hàng
        \textbf{mọi lúc}, không xét gì thêm. \note{(UAT: đang BẬT)}
  \item \textbf{Khung giờ theo khu vực}. Nếu khu vực của cửa hàng có ít nhất một khung
        (\rt{wujia.order.window}) đang hoạt động, hệ thống chỉ xét các khung đó và cho đặt khi
        \textbf{bất kỳ khung nào} đang mở. Một khu vực được phép có nhiều khung (sáng + chiều).
  \item \textbf{Khung giờ chung}. Khu vực không có khung riêng thì dùng cặp tham số chung
        \rt{portal_order_time_from} / \rt{portal_order_time_to}. \note{(UAT: 10:00--04:00)}
\end{enumerate}

Khung có giờ kết thúc \textbf{nhỏ hơn} giờ bắt đầu được hiểu là \textbf{vắt qua nửa đêm} (ví dụ
10:00 hôm nay đến 04:00 ngày mai); màn hiển thị ghi rõ ``hôm nay'' / ``ngày mai'' kèm ``UTC+7''
để không hiểu nhầm.

\textbf{Hiện trạng trên UAT:}

\begin{center}\footnotesize
\begin{tabular}{L{3.9cm}L{3.3cm}L{3.9cm}L{3.5cm}}
\rowcolor{hdrbg}\textbf{Cửa hàng} & \textbf{Khu vực} & \textbf{Khung giờ áp dụng} &
  \textbf{Nguồn}\\ \midrule
TP HCM Quận 1 (HCM-01) & TP HCM Quận 1--3 & 00:00 -- 23:50 & Khung của khu vực \\
Cầu Giấy (HN-01) & Hà Nội Trung tâm & 10:00 -- 04:00 & Khung chung \\
Đống Đa (HN-02) & Hà Nội Trung tâm & 10:00 -- 04:00 & Khung chung \\
Cửa hàng TEST Issue 134 & Hà Nội Trung tâm & 10:00 -- 04:00 & Khung chung \\
\end{tabular}
\end{center}

\mhoi

\begin{ask}
\begin{enumerate}[leftmargin=*]
  \item \textbf{Tên khung giờ không khớp giờ thật.} Khung duy nhất trên UAT \textbf{tên là}
        ``00:00 - 05:00'' nhưng giờ thật đang là \textbf{00:00 -- 23:50}. Tên là chữ người dùng
        tự gõ, hệ thống không kiểm --- và chính chuỗi tên này hiện lên portal. Có cần bắt hệ
        thống tự sinh tên theo giờ không?
  \item \textbf{Giờ được tính theo múi giờ của TÀI KHOẢN đăng nhập}, không theo khu vực của cửa
        hàng, trong khi dải thông báo luôn ghi ``UTC+7''. Tài khoản chưa đặt múi giờ sẽ bị tính
        theo giờ UTC --- lệch 7 tiếng so với chữ hiển thị. Đề nghị chốt: lấy múi giờ theo
        \textbf{khu vực cửa hàng}, hay ép mọi tài khoản portal về giờ Việt Nam?
  \item \textbf{Khu vực Hà Nội (3/4 cửa hàng) chưa có khung riêng} nên đang chạy khung chung. BA
        xác nhận đây là ý muốn, hay cần cấu hình khung riêng cho Hà Nội?
\end{enumerate}
\end{ask}
